\documentclass[lettersize,journal]{IEEEtran}

\usepackage{amsmath,amsfonts,amssymb}
\usepackage{algorithmic}
\usepackage{algorithm}
\usepackage{array}
\usepackage[caption=false,font=normalsize,labelfont=sf,textfont=sf]{subfig}
\usepackage{textcomp}
\usepackage{stfloats}
\usepackage{url}
\usepackage{graphicx}
\usepackage{cite}
\usepackage{bm}
\usepackage{booktabs}

\hyphenation{op-tical net-works semi-conduc-tor}

\begin{document}

\title{DCPAM: A Dual-Camera-Based Point-to-Axis Measurement Model
       with Laser Reference Line}

% TODO: 补充作者信息、单位、邮箱
\author{
  Author~One,
  Author~Two,
  and~Author~Three
  \thanks{Manuscript received XXXX XX, 2025.}
  \thanks{The authors are with XXXX University, China.
          (e-mail: xxxx@xxxx.edu.cn)}
}

\markboth{Journal of XXX,~Vol.~X, No.~X, 2025}%
{Author \MakeLowercase{\textit{et al.}}:
 DCPAM: Dual-Camera Point-to-Axis Measurement}

\maketitle

% ============================================================
%  Abstract & Keywords
% ============================================================
\begin{abstract}
  In large-scale rotating machinery such as turbines, generators,
  and compressors, shaft alignment accuracy is critical to
  operational balance and equipment lifespan.
  %
  Traditional steel-wire alignment methods suffer from
  wire sag, vibration sensitivity, and limited measurement range.
  %
  This paper proposes DCPAM
  (Dual-Camera-based Point-to-Axis Measurement),
  a non-contact measurement model that replaces the steel wire
  with a laser reference line and employs a dual-camera system
  to reconstruct the three-dimensional laser axis.
  %
  The model computes the perpendicular distance from any target
  point to the reconstructed laser axis through a pipeline of
  pixel extraction, camera-to-world coordinate transformation,
  mirror-reflection virtual-image mapping, and cross-product
  distance calculation.
  %
  Two parameter-tuning strategies are presented:
  DCPAM-CV, an experience-driven classical computer-vision
  workflow, and DCPAM-CNN, a convolutional-neural-network-based
  approach.
  %
  A closed-form error propagation analysis and Monte Carlo
  validation quantify the measurement uncertainty under
  micrometer-level perturbations.
  % TODO: 补充实验结果摘要——最终精度数据
\end{abstract}

\begin{IEEEkeywords}
  Shaft alignment, laser measurement, dual-camera system,
  point-to-axis distance, error propagation, computer vision.
\end{IEEEkeywords}

% ============================================================
%  I. Introduction
% ============================================================
\section{Introduction}

\IEEEPARstart{I}{n} the power generation, marine, and aerospace
industries, large rotating equipment---turbines, generators,
compressors---features long shafts whose installation accuracy
directly governs operational balance and service life.
%
Shaft alignment (or ``centreline finding'') is therefore a
critical step during installation and commissioning of such
machinery~\cite{piotrowski2006shaft}.

The prevailing industrial practice remains the
\emph{steel-wire method}, in which a taut piano wire is
stretched along the nominal shaft axis and the radial distance
from each bearing journal to the wire is measured with a
micrometer or dial indicator.
%
Despite its simplicity, the steel-wire method has well-known
drawbacks:
gravitational sag introduces systematic curvature
that must be compensated;
the wire is sensitive to air currents and mechanical vibration;
and the achievable measurement range is limited by wire tension
and sag correction accuracy.

Laser-based alignment offers a promising alternative.
%
By replacing the steel wire with a laser beam as the reference
axis, the problems of sag and vibration sensitivity are
eliminated.
%
However, the laser beam itself is invisible in free space;
its position can only be observed where it strikes a
receiving surface.
%
This paper proposes \textbf{DCPAM}
(Dual-Camera-based Point-to-Axis Measurement),
a model that employs two cameras---each observing the laser
spot on its own receiving screen---to reconstruct the
three-dimensional laser axis and compute the perpendicular
distance from an arbitrary target point to that axis.

The contributions of this work are as follows:
%
(1)~A complete mathematical formulation of the dual-camera
laser measurement pipeline, from pixel extraction to
point-to-axis distance computation.
%
(2)~Two parameter-tuning strategies:
DCPAM-CV (classical computer vision) and
DCPAM-CNN (convolutional neural network).
%
(3)~A closed-form error propagation analysis with
Monte Carlo validation.
%
(4)~A tilt-deviation rejection model that determines
the maximum allowable probe inclination for a given
target accuracy.

% TODO: 补充文献综述——现有激光对中方法、双目视觉测量、工业对中标准

% ============================================================
%  II. DCPAM Model Overview
% ============================================================
\section{DCPAM Model Overview}

The DCPAM model takes as input the images from two cameras and
the length of a measurement tool bar, and outputs the
perpendicular distance from the target point to the laser
reference axis:
%
\begin{equation}
  \boxed{h_t = \mathrm{DCPAM}(I_{C_1},\, I_{C_2},\, l_T)}
  \label{eq:dcpam-overview}
\end{equation}
%
where $h_t$ is the target-point-to-axis distance,
$I_{C_1}$ and $I_{C_2}$ are the images captured by cameras
$C_1$ and $C_2$ respectively,
and $l_T$ is the length of the tool bar (side gauge).

Two parameter-tuning strategies are considered:

\textbf{DCPAM-CV.}
An experience-driven workflow in which the operator
iteratively adjusts CV algorithm parameters
(thresholds, filter sizes, fitting methods),
evaluates on a test set, and manually refines the
pipeline until the target accuracy is achieved.

\textbf{DCPAM-CNN.}
A deep-learning approach that trains a convolutional
neural network to directly regress the laser spot
centre from raw images, replacing hand-tuned CV
processing.

% TODO: 补充 DCPAM 整体流程图（建议画一个 pipeline figure）

% ============================================================
%  III. DCPAM-CV: Classical Vision Pipeline
% ============================================================
\section{DCPAM-CV: Classical Vision Pipeline}

This section details the complete measurement pipeline
of the classical computer-vision variant.
%
Fig.~\ref{fig:pipeline} illustrates the data flow from raw
images to the final distance measurement.

% TODO: 插入 pipeline 流程图
% \begin{figure}[!t]
%   \centering
%   \includegraphics[width=\columnwidth]{pipeline}
%   \caption{DCPAM-CV measurement pipeline.}
%   \label{fig:pipeline}
% \end{figure}


% ------------------------------------------------------------
\subsection{Coordinate System Definitions}
\label{sec:coord}

We define three coordinate frames, all right-handed:

\textbf{Camera coordinate frame $C1$.}
Origin at the optical centre of camera $C_1$;
$Z$-axis aligned with the optical axis.
%
A point $P$ in this frame is denoted
$P^{C1} = (x_1, y_1, z_1)^T$.

\textbf{Camera coordinate frame $C2$.}
Origin at the optical centre of camera $C_2$;
$Z$-axis aligned with its optical axis.
%
A point in this frame is denoted
$P^{C2} = (x_2, y_2, z_2)^T$.

\textbf{Device coordinate frame $D$.}
By convention, $D \equiv C1$---the device frame coincides
with the front camera coordinate frame.
%
All final computations are performed in $D$.

The optical centres of the two cameras in their own frames are
trivially
$O_{C_1} = (0,0,0)^T$ and $O_{C_2} = (0,0,0)^T$.
%
Note that this is a definitional identity:
the intrinsic parameters (focal length, principal point,
distortion) govern the projection relationship but do not
affect the location of the optical centre in the camera's
own coordinate frame.


% ------------------------------------------------------------
\subsection{Key Points}
\label{sec:keypoints}

We define the following key points (all expressed in $C1$
unless otherwise noted):

$P^f_{\mathrm{real}}$:
the laser spot on the front receiving screen
(real image in $C1$).

$P^f_{\mathrm{virtual}}$:
the virtual-image point obtained by rotating
$P^f_{\mathrm{real}}$ about the mirror hinge
(virtual image in $C1$).

$P^{b\text{-}C2}_{\mathrm{real}}$:
the laser spot on the rear receiving screen
(real image in $C2$).

$P^{b\text{-}C2}_{\mathrm{virtual}}$:
the virtual-image point of the rear screen
(virtual image in $C2$).

$P^b_{\mathrm{virtual}}$:
$P^{b\text{-}C2}_{\mathrm{virtual}}$ transformed
into $C1$.

$P_T$:
the target point whose distance to the laser axis
is to be measured.


% ------------------------------------------------------------
\subsection{Laser Spot Centre Extraction}
\label{sec:extraction}

Let the pixel coordinates of the laser spot in
$C_1$ and $C_2$ be $(u_1, v_1)$ and $(u_2, v_2)$,
with homogeneous representations
%
\begin{equation}
  \mathbf{p}_1 =
  \begin{bmatrix} u_1 \\ v_1 \\ 1 \end{bmatrix}, \quad
  \mathbf{p}_2 =
  \begin{bmatrix} u_2 \\ v_2 \\ 1 \end{bmatrix}.
  \label{eq:pixel-coords}
\end{equation}

% TODO: 补充光斑提取的具体 CV 算法
%       （高斯拟合、质心法、Canny 边缘 + 梯度投票等，
%         参考 center.py 中的实现）


% ------------------------------------------------------------
\subsection{Pixel-to-Camera Coordinate Transformation}
\label{sec:pixel2cam}

Let $K_1$ and $K_2$ be the intrinsic matrices of cameras
$C_1$ and $C_2$, and let $Z_s$ be the known distance from
the camera optical centre to the receiving screen along the
optical axis.
%
The back-projection from pixel coordinates to camera
coordinates is
%
\begin{equation}
  \boxed{
    P^f_{\mathrm{real}} = Z_s \, K_1^{-1}\, \mathbf{p}_1, \quad
    P^{b\text{-}C2}_{\mathrm{real}}
      = Z_s \, K_2^{-1}\, \mathbf{p}_2
  }
  \label{eq:backproj}
\end{equation}

\textbf{Dependencies.}
Camera calibration provides $K_1$, $K_2$;
the tool geometry provides $Z_s$.

\textbf{Error sources.}
The accuracy of $Z_s$ depends on the mechanical precision
of the probe fixture.


% ------------------------------------------------------------
\subsection{Real-to-Virtual Point Transformation}
\label{sec:real2virtual}

The receiving screen is mounted on a mirror that can rotate
about a hinge axis.
%
To reconstruct the ``unfolded'' laser path,
each real image point must be rotated to its virtual position.

In frame $C1$, let the rotation centre be
$R_{O1} = (x_{O1}, 0, z_{O1})^T$,
and let $\theta$ be the rotation angle
(counter-clockwise from $Z$ towards $X$ in a right-handed
frame).
%
Then
%
\begin{equation}
  \boxed{P^f_{\mathrm{virtual}} = H_1 \cdot P^f_{\mathrm{real}}}
  \label{eq:H1}
\end{equation}
%
where $H_1$ is a $4\times4$ rigid-body transformation matrix
combining rotation and translation:
%
\begin{equation}
  H_1 =
  \begin{bmatrix}
    \cos\theta  & 0 & \sin\theta
      & x_{O1}(1{-}\cos\theta) - z_{O1}\sin\theta \\
    0           & 1 & 0           & 0 \\
    -\sin\theta & 0 & \cos\theta
      & z_{O1}(1{-}\cos\theta) + x_{O1}\sin\theta \\
    0           & 0 & 0           & 1
  \end{bmatrix}
  \label{eq:H1-matrix}
\end{equation}

The transformation represents a rotation about an axis
parallel to $Y$ that passes through $R_{O1}$.

Similarly, the rear real-image point is rotated in $C2$
about the centre $R_{O2} = (x_{O2}, 0, z_{O2})^T$
by angle $2\theta$:
%
\begin{equation}
  P^{b\text{-}C2}_{\mathrm{virtual}}
    = H_2 \cdot P^{b\text{-}C2}_{\mathrm{real}}
  \label{eq:H2}
\end{equation}

\textbf{Dependencies.}
Geometric dimensions $R_{O1}$, $R_{O2}$, and the
rotation angle $\theta$.

\textbf{Ideal-case simplification.}
Under ideal assembly,
$z_{O1} = z_{O2} = Z_s$.

% TODO: 补充旋转推导过程的详细数学证明（可放入附录）


% ------------------------------------------------------------
\subsection{$C2$-to-$C1$ Coordinate Transformation}
\label{sec:c2toc1}

The virtual-image point $P^{b\text{-}C2}_{\mathrm{virtual}}$
is expressed in frame $C2$ and must be converted to $C1$:
%
\begin{equation}
  \boxed{
    P^b_{\mathrm{virtual}}
      = T_{C2 \to C1}
        \cdot P^{b\text{-}C2}_{\mathrm{virtual}}
  }
  \label{eq:T-c2c1}
\end{equation}

The transformation matrix is computed from the extrinsic
parameters of the two cameras with respect to a common world
frame $W$:
%
\begin{equation}
  T_{C2 \to C1} = E_1 \cdot E_2^{-1}
  \label{eq:T-from-ext}
\end{equation}
%
where $E_1$, $E_2$ are the $4\times4$ extrinsic matrices
(world-to-camera) of $C_1$ and $C_2$ respectively.

Equivalently, if the relative pose is parameterised as
rotation $R_{21}$ and translation $t_{21}$
(from $C1$ to $C2$), then
%
\begin{align}
  R_{12} &= R_{21}^T  \label{eq:R12}\\
  t_{12} &= -R_{21}^T \, t_{21}  \label{eq:t12}
\end{align}
%
and $P^{C1} = R_{12}\, P^{C2} + t_{12}$.

\textbf{Dependencies.}
Stereo camera calibration (extrinsic parameters).


% ------------------------------------------------------------
\subsection{Point-to-Axis Distance Computation}
\label{sec:distance}

With $P^f_{\mathrm{virtual}}$, $P^b_{\mathrm{virtual}}$,
and $P_T$ all expressed in $C1$,
the perpendicular distance from $P_T$ to the laser axis
(the line through
$P^f_{\mathrm{virtual}}$ and $P^b_{\mathrm{virtual}}$)
is given by the classical cross-product formula:
%
\begin{equation}
  \boxed{
    H = \frac{
      \left\| \vec{L} \times \vec{w} \right\|
    }{
      \left\| \vec{L} \right\|
    }
  }
  \label{eq:distance}
\end{equation}
%
where
%
\begin{align}
  \vec{L} &= P^f_{\mathrm{virtual}} - P^b_{\mathrm{virtual}}
  \label{eq:L-vec} \\
  \vec{w} &= P_T - P^b_{\mathrm{virtual}}
  \label{eq:w-vec}
\end{align}

The vector $\vec{L}$ is the direction vector of the
reconstructed laser axis in frame $C1$.

\textbf{Dependencies.}
The target point $P_T$ is determined by the tool-bar
geometry and the measurement configuration.

% ============================================================
%  IV. Precision Analysis
% ============================================================
\section{Precision Analysis}
\label{sec:precision}

% TODO: 补充精度分析的整体流程图

\subsection{End-to-End Matrix Formulation}
\label{sec:matrix-form}

Combining the transformations from
Sections~\ref{sec:pixel2cam}--\ref{sec:c2toc1},
the two virtual points can be expressed as linear functions
of the pixel coordinates:
%
\begin{align}
  P^f_{\mathrm{virtual}}
    &= H_1
       \begin{bmatrix}
         Z_s\, K_1^{-1}\, \mathbf{p}_1 \\ 1
       \end{bmatrix}
  \label{eq:pf-end2end} \\[6pt]
  P^b_{\mathrm{virtual}}
    &= T_{C2 \to C1} \cdot H_2
       \begin{bmatrix}
         Z_s\, K_2^{-1}\, \mathbf{p}_2 \\ 1
       \end{bmatrix}
  \label{eq:pb-end2end}
\end{align}

Since $Z_s K_i^{-1}$ and the rigid-body matrices are all
constant for a given configuration,
each coordinate of
$P^f_{\mathrm{virtual}} = (x_f, y_f, z_f)^T$ and
$P^b_{\mathrm{virtual}} = (x_b, y_b, z_b)^T$
is an \emph{affine} function of the pixel coordinates:
%
\begin{align}
  x_f &= A_1\, u_1 + B_1\, v_1 + C_1 \notag\\
  y_f &= D_1\, u_1 + E_1\, v_1 + F_1 \notag\\
  z_f &= G_1\, u_1 + H_1\, v_1 + I_1
  \label{eq:affine-f}
\end{align}
\begin{align}
  x_b &= A_2\, u_2 + B_2\, v_2 + C_2 \notag\\
  y_b &= D_2\, u_2 + E_2\, v_2 + F_2 \notag\\
  z_b &= G_2\, u_2 + H_2\, v_2 + I_2
  \label{eq:affine-b}
\end{align}
%
where the coefficients $A_1, B_1, \ldots, I_2$ are
determined by the composite matrix products.

The target point is
%
\begin{equation}
  P_T =
  \begin{bmatrix} x_0 \\ y_0 \\ z_0 + l \end{bmatrix}
  \label{eq:pt}
\end{equation}
%
where $(x_0, y_0, z_0)^T$ is the tool-bar mounting position
and $l = l_T$ is the tool-bar length.


\subsection{Expanded Distance Expression}
\label{sec:expanded}

Substituting the affine expressions into
\eqref{eq:distance}--\eqref{eq:w-vec},
the measurement $H$ becomes an explicit function of
five variables:
%
\begin{equation}
  H(u_1, v_1, u_2, v_2, l)
    = \frac{
        \left\| \vec{d} \times \vec{t} \right\|
      }{
        \|\vec{d}\|
      }
  \label{eq:H-expanded}
\end{equation}
%
with
%
\begin{equation}
  \vec{d} =
  \begin{bmatrix}
    (A_2 u_2 - A_1 u_1)+(B_2 v_2 - B_1 v_1)+(C_2-C_1) \\
    (D_2 u_2 - D_1 u_1)+(E_2 v_2 - E_1 v_1)+(F_2-F_1) \\
    (G_2 u_2 - G_1 u_1)+(H_2 v_2 - H_1 v_1)+(I_2-I_1)
  \end{bmatrix}
  \label{eq:d-vec}
\end{equation}
%
\begin{equation}
  \vec{t} =
  \begin{bmatrix}
    -A_1 u_1 - B_1 v_1 + (x_0 - C_1) \\
    -D_1 u_1 - E_1 v_1 + (y_0 - F_1) \\
    -G_1 u_1 - H_1 v_1 + l + (z_0 - I_1)
  \end{bmatrix}
  \label{eq:t-vec}
\end{equation}

\textbf{Key observation.}
When $l$ is fixed and only the pixel coordinates
$u_1, v_1, u_2, v_2$ vary
(e.g.\ due to pixel extraction noise),
the resulting change $\Delta H$ is
\emph{essentially independent of $l$}.
%
The absolute value of $H$ depends on $l$,
but the measurement uncertainty $\sigma_H$
depends only on the pixel noise magnitude and the
system coefficients $A_1, \ldots, I_2$.

% TODO: 用具体数值做一个数值示例证明以上结论


\subsection{Error Propagation}
\label{sec:error-prop}

By first-order Taylor expansion, the variance of $H$ is
%
\begin{equation}
  \sigma_H^2 \approx
    \sum_{q \in \mathcal{Q}}
    \left(\frac{\partial H}{\partial q}\right)^2
    \sigma_q^2
  \label{eq:error-prop}
\end{equation}
%
where
$\mathcal{Q} = \{u_1, v_1, u_2, v_2\}$
(or equivalently $\{y_f, z_f, y_b, z_b\}$ after
coordinate transformation),
and $\sigma_q$ is the standard deviation of each
input parameter.

The partial derivatives can be computed analytically
from~\eqref{eq:H-expanded} or numerically via finite
differences.

% TODO: 补充解析偏导数的显式表达式


\subsection{Monte Carlo Validation}
\label{sec:mc}

To validate the linear error propagation, we perform
a Monte Carlo simulation:
%
$N = 10{,}000$ samples are drawn with each pixel coordinate
perturbed uniformly within $\pm 1\,\mu\mathrm{m}$
(equivalent pixel noise).
%
The empirical standard deviation $\hat{\sigma}_H$ is compared
with the analytical $\sigma_H$ from~\eqref{eq:error-prop}.

% TODO: 补充 Monte Carlo 结果表格和直方图
%       （参考 scripts/run_analysis.py 的输出）


\subsection{Sensitivity Analysis}
\label{sec:sensitivity}

Each input parameter is perturbed by $\pm 1\,\mu\mathrm{m}$
independently, and the resulting $|\Delta H|$ is recorded.
%
The variance contribution of each parameter is
%
\begin{equation}
  \eta_q
    = \frac{
        \left(\frac{\partial H}{\partial q}\right)^2 \sigma_q^2
      }{
        \sigma_H^2
      }
    \times 100\%
  \label{eq:variance-contrib}
\end{equation}

% TODO: 补充灵敏度分析的实验数据表格和图表

% ============================================================
%  V. Camera Calibration
% ============================================================
\section{Camera Calibration}
\label{sec:calibration}

Accurate camera calibration is a prerequisite for the DCPAM
pipeline.
%
This section describes the calibration of intrinsic
parameters, extrinsic parameters, and the screen-plane
distance $Z_s$.


\subsection{Intrinsic Parameter Calibration}
\label{sec:intrinsic}

Each camera is modelled with the SIMPLE\_RADIAL model:
focal length $f$, principal point $(c_x, c_y)$,
and a single radial distortion coefficient $k$.
%
The intrinsic matrix is
%
\begin{equation}
  K =
  \begin{bmatrix}
    f & 0 & c_x \\
    0 & f & c_y \\
    0 & 0 & 1
  \end{bmatrix}
  \label{eq:K}
\end{equation}

Calibration is performed using COLMAP's feature extraction
and Structure-from-Motion pipeline on a set of checkerboard
images captured from multiple viewpoints.

% TODO: 补充标定板规格、标定图片数量、重投影误差等实验细节
% TODO: 补充畸变矫正的具体处理方法


\subsection{Extrinsic Parameter Calibration}
\label{sec:extrinsic}

The extrinsic parameters describe the pose of each camera
relative to a common world frame $W$:
%
\begin{equation}
  E_i =
  \begin{bmatrix}
    R_i & t_i \\
    \mathbf{0}^T & 1
  \end{bmatrix},
  \quad i = 1, 2
  \label{eq:extrinsic}
\end{equation}

The relative transformation $T_{C2 \to C1}$ is then obtained
via~\eqref{eq:T-from-ext}.

The dual-camera calibration procedure follows these steps:

(1)~Capture synchronised image pairs from both cameras.

(2)~Extract and match features across all images simultaneously.

(3)~Run SfM mapping with undistortion.

(4)~Export intrinsic and extrinsic parameters.

(5)~Compute relative pose (expected baseline $\sim$7.5--8.0\,cm).

% TODO: 补充 COLMAP 标定的具体步骤截图和参数设置
% TODO: 补充标定精度验证实验


\subsection{Screen-Plane Distance Calibration}
\label{sec:zs}

The distance $Z_s$ from the camera optical centre to the
receiving screen along the optical axis is a critical
parameter, as it scales the back-projection~\eqref{eq:backproj}.
%
$Z_s$ is determined by the mechanical dimensions of the
probe fixture and verified by measurement.

The composite matrix $Z_s\, K_i^{-1}$ can also be calibrated
directly by observing known points on the receiving screen.

% TODO: 补充 Zs 标定方法的详细描述和精度分析


\subsection{Camera Tilt Angle Calibration}
\label{sec:tilt-calib}

When the camera is mounted in a nominally horizontal
orientation, a residual tilt angle $\alpha$ may exist
between the camera's $x$-pixel direction and the
horizontal direction of the world frame.
%
This angle must be calibrated and compensated.

The relationship between a world-frame point $(X, Y)$ and
the camera-frame point $(x, y)$ under a planar rotation is
%
\begin{equation}
  \begin{bmatrix} x \\ y \end{bmatrix}
  =
  \begin{bmatrix}
    \cos\alpha & -\sin\alpha \\
    \sin\alpha &  \cos\alpha
  \end{bmatrix}
  \begin{bmatrix} X \\ Y \end{bmatrix}
  \label{eq:tilt-rotation}
\end{equation}

% TODO: 补充 alpha 角度标定的实验方法

% ============================================================
%  VI. Mirror Transformation Calibration
% ============================================================
\section{Mirror Transformation Calibration}
\label{sec:mirror-calib}

The matrices $H_1$ and $H_2$
(Section~\ref{sec:real2virtual}) encode the mirror
rotation.
%
Their calibration requires knowledge of:
the rotation centre position $R_{O1}$
(or $R_{O2}$);
the rotation angle $\theta$.

% TODO: 补充 H1, H2 标定方法
% TODO: 补充旋转中心和旋转角度的测量方法

% ============================================================
%  VII. DCPAM-CNN
% ============================================================
\section{DCPAM-CNN: Deep-Learning-Based Approach}
\label{sec:cnn}

While DCPAM-CV relies on hand-crafted feature extraction
and manual parameter tuning,
DCPAM-CNN trains a convolutional neural network to directly
regress the laser spot centre coordinates from raw camera
images.

% TODO: 补充以下内容：
%   - CNN 网络结构（输入、卷积层、全连接层、输出）
%   - 损失函数设计
%   - 训练数据集构建方法
%   - 训练超参数
%   - 与 DCPAM-CV 的精度对比实验

% ============================================================
%  VIII. Tilt Deviation Model
% ============================================================
\section{Tilt Deviation and Measurement Accuracy}
\label{sec:tilt}

In practice, the measurement probe may not be perfectly
aligned with the nominal direction.
%
A tilt deviation $\beta$ introduces systematic error into
the distance measurement.
%
This section derives the relationship between $\beta$ and
the measurement accuracy, and determines the maximum
allowable tilt for a specified accuracy target $\sigma$.

Given a target accuracy $\sigma$,
the \emph{tilt rejection threshold} $\beta(\sigma)$ is
defined as the maximum inclination angle beyond which
the measured data should be discarded to maintain the
required precision.

% TODO: 补充以下内容：
%   - 倾斜角 β 的数学定义和几何示意图
%   - β 对测量误差 H 的影响推导
%   - β(σ) 阈值的解析表达式或数值求解方法
%   - 不同 σ 目标下的 β 阈值表格

% ============================================================
%  IX. Experiments
% ============================================================
\section{Experiments}
\label{sec:experiments}

% TODO: 补充以下实验内容：
%   1. 实验平台描述（相机型号、激光器、工装尺寸）
%   2. 标定实验结果（内参、外参、Zs）
%   3. 静态精度实验（已知距离 vs 测量距离）
%   4. 重复性实验（多次测量同一点的标准差）
%   5. 动态测量实验（不同位置的测量结果）
%   6. DCPAM-CV vs DCPAM-CNN 对比
%   7. 与传统钢丝法的对比实验


\subsection{Experimental Setup}

The experimental platform consists of a Daheng Galaxy
industrial camera pair (resolution $2592 \times 1944$),
a collimated laser source, receiving screens with
mirror mounts, and a precision linear stage for
ground-truth positioning.

% TODO: 补充实验装置照片和示意图


\subsection{Calibration Results}

Table~\ref{tab:calib} summarises the calibration results
for both cameras.

\begin{table}[!t]
  \caption{Camera Calibration Parameters
           \label{tab:calib}}
  \centering
  \small
  \begin{tabular}{lcc}
    \toprule
    Parameter & $C_1$ (Left) & $C_2$ (Right) \\
    \midrule
    Model            & SIMPLE\_RADIAL & SIMPLE\_RADIAL \\
    Resolution (px)  & $2592\times1944$ & $2592\times1944$ \\
    Focal length (px) & 2918.30 & 3757.72 \\
    Principal point  & (1296, 972) & (1296, 972) \\
    Distortion $k$   & $-0.2137$ & $-0.2384$ \\
    \bottomrule
  \end{tabular}
\end{table}

% TODO: 补充外参标定结果（旋转矩阵、平移向量、基线长度）
% TODO: 补充标定重投影误差


\subsection{Measurement Accuracy}

% TODO: 补充精度实验结果
%   - 测量值 vs 真值对比表
%   - 误差统计（均值、标准差、最大误差）
%   - 误差分布直方图

% ============================================================
%  X. Conclusion
% ============================================================
\section{Conclusion}

This paper presented DCPAM, a dual-camera-based
point-to-axis measurement model that replaces the
traditional steel wire with a laser reference line.
%
The complete mathematical pipeline---from pixel
extraction through coordinate transformation to
point-to-line distance computation---was derived
and validated through error propagation analysis
and Monte Carlo simulation.

% TODO: 补充结论——总结实验结果、最终达到的精度、
%       与传统方法的对比优势、未来工作方向


% ============================================================
%  Acknowledgments
% ============================================================
\section*{Acknowledgments}

% TODO: 补充致谢信息


% ============================================================
%  References
% ============================================================
\begin{thebibliography}{10}
\bibliographystyle{IEEEtran}

\bibitem{piotrowski2006shaft}
J.~Piotrowski,
\textit{Shaft Alignment Handbook},
3rd~ed.\hskip 1em
Boca Raton, FL, USA: CRC Press, 2006.

% TODO: 补充更多参考文献：
%   - 激光对中相关文献
%   - 双目视觉测量文献
%   - 相机标定文献（Zhang's method, COLMAP）
%   - CNN 回归相关文献
%   - 误差传播理论文献
%   - 工业对中标准（ISO, GB）

\end{thebibliography}


\end{document}
